\documentclass[12pt,a4paper]{article}

% ============================================================
% PACKAGES
% ============================================================
\usepackage[utf8]{inputenc}
\usepackage[T5]{fontenc}
\usepackage[vietnamese]{babel}
\usepackage{geometry}
\usepackage{amsmath}
\usepackage{amssymb}
\usepackage{booktabs}
\usepackage{array}
\usepackage{longtable}
\usepackage{hyperref}
\usepackage{enumitem}
\usepackage{setspace}
\usepackage{parskip}
\usepackage{titlesec}
\usepackage{xcolor}
\usepackage{microtype}

% ============================================================
% PAGE GEOMETRY
% ============================================================
\geometry{
  a4paper,
  top=2.5cm,
  bottom=2.5cm,
  left=3cm,
  right=2.5cm
}

% ============================================================
% HYPERREF SETUP
% ============================================================
\hypersetup{
  colorlinks=true,
  linkcolor=blue!70!black,
  urlcolor=blue!70!black,
  citecolor=blue!70!black,
  pdftitle={Cơ sở lý thuyết},
  pdfauthor={},
  pdfencoding=auto
}

% ============================================================
% SPACING
% ============================================================
\onehalfspacing
\setlength{\parindent}{0pt}
\setlength{\parskip}{6pt}

% ============================================================
% SECTION FORMATTING
% ============================================================
\titleformat{\section}{\large\bfseries}{\thesection}{1em}{}
\titleformat{\subsection}{\normalsize\bfseries}{\thesubsection}{1em}{}
\titleformat{\subsubsection}{\normalsize\bfseries}{\thesubsubsection}{1em}{}

% ============================================================
% TABLE COLUMN TYPES
% ============================================================
\newcolumntype{L}[1]{>{\raggedright\arraybackslash}p{#1}}
\newcolumntype{C}[1]{>{\centering\arraybackslash}p{#1}}

% ============================================================
% DOCUMENT
% ============================================================
\begin{document}

% ============================================================
\section{Cơ sở lý thuyết}
% ============================================================

% ------------------------------------------------------------
\subsection{Giới thiệu bộ dữ liệu}
% ------------------------------------------------------------

Báo cáo sử dụng tệp \texttt{All\_GPUs.csv} được trích xuất từ bộ dữ liệu
\textit{``Computer Parts (CPUs and GPUs)''} trên nền tảng Kaggle (tác giả: Iliasse Kkaf).

Tập dữ liệu cung cấp các thông số cấu hình thực tế của card đồ họa (GPU) từ bốn nhà sản
xuất: NVIDIA, AMD, Intel và ATI~-- hãng ATI sau đó được gộp chung vào AMD trong quá trình
làm sạch dữ liệu. Các trường đặc trưng chính phục vụ cho phân tích bao gồm: tốc độ xung
nhịp lõi, công suất tiêu thụ (TDP), thông số bộ nhớ (dung lượng, chuẩn, băng thông), thời
điểm ra mắt và giá phát hành. Do đặc thù dữ liệu thô chứa nhiều giá trị khuyết, tập dữ
liệu sẽ trải qua quy trình tiền xử lý nghiêm ngặt trước khi được đưa vào các mô hình
thống kê.

% ------------------------------------------------------------
\subsection{Lựa chọn biến phân tích và giải thích}
% ------------------------------------------------------------

\textbf{Câu hỏi nghiên cứu:}

\begin{itemize}[leftmargin=1.5em]
  \item Những thông số kỹ thuật nào (xung nhịp, băng thông bộ nhớ, công suất, năm phát
        hành, hãng sản xuất, loại bộ nhớ) thực sự ảnh hưởng đến giá phát hành GPU?

  \item Có sự khác biệt có ý nghĩa thống kê về giá phát hành trung bình giữa hai hãng
        NVIDIA và AMD hay không?
\end{itemize}

Để trả lời các câu hỏi này, nhóm lựa chọn 7 biến độc lập và 1 biến phụ thuộc dựa trên
hai tiêu chí: (1)~có liên quan trực tiếp đến giá phát hành theo lý thuyết kinh tế học
phần cứng, và (2)~có đủ dữ liệu không bị khuyết để xây dựng mô hình hồi quy ổn định.
Bảng~\ref{tab:bien} sau trình bày chi tiết từng biến.

\begin{table}[ht]
\centering
\caption{Danh sách biến sử dụng trong mô hình và vai trò tương ứng}
\label{tab:bien}
\renewcommand{\arraystretch}{1.3}
\begin{tabular}{L{2.4cm} L{2.8cm} L{5.2cm} L{3.6cm}}
\toprule
\textbf{Biến} & \textbf{Kiểu thực tế} & \textbf{Lý do / Giải thích} & \textbf{Vai trò trong mô hình} \\
\midrule
\texttt{manufacturer}  & Định danh (Nominal)  
  & Phân nhóm hãng sản xuất (NVIDIA, AMD, Intel)  
  & Biến độc lập (phân loại) \\

\texttt{memory\_type}  & Định danh (Nominal)  
  & Các chuẩn VRAM (GDDR5, GDDR3, DDR3, HBM\ldots{} và các chuẩn khác) -- số lượng giá trị
    phân biệt có giới hạn.  
  & Biến độc lập (phân loại), đã được xem xét và loại bỏ do VIF = 8.12; xem Phần~4, mục 7.1.2 \\

\texttt{memory\_bus}   & Rời rạc (Discrete)   
  & Băng thông bộ nhớ (bit) -- giá trị rời rạc tuân theo các chuẩn kỹ thuật của ngành.  
  & Biến độc lập (số) \\

\texttt{memory\_size}  & Rời rạc (Discrete)   
  & Dung lượng VRAM (GB hoặc MB) -- giá trị rời rạc phân bố theo thiết kế của nhà sản xuất.  
  & Biến độc lập (số) \\

\texttt{tdp}           & Liên tục (Continuous)
  & Công suất tiêu thụ điện (Watt) -- ảnh hưởng trực tiếp đến chi phí sản xuất  
  & Biến độc lập (số) \\

\texttt{core\_speed}   & Liên tục (Continuous)
  & Tốc độ xung nhịp lõi (MHz) -- phản ánh hiệu năng xử lý  
  & Biến độc lập (số) \\

\texttt{release\_year} & Rời rạc $\to$ Số     
  & Ngày phát hành được parse thành năm (integer) để nắm xu hướng thời gian  
  & Biến độc lập (số) \\

\texttt{release\_price}& Liên tục (Continuous)
  & Giá phát hành (USD) -- biến mục tiêu cần giải thích  
  & Biến phụ thuộc $Y$ \\
\bottomrule
\end{tabular}
\end{table}

\textbf{\textit{Lưu ý đặc biệt:}}

Biến \texttt{memory\_bus} tuy mang giá trị số nhưng thực chất là biến rời rạc vì chỉ
tồn tại một số lượng giá trị cố định theo chuẩn kỹ thuật của ngành. Trong mô hình hồi
quy, biến này được giữ nguyên dạng số vì thứ tự và khoảng cách giữa các giá trị có ý
nghĩa vật lý thực sự (băng thông tăng gấp đôi phản ánh hiệu năng tăng tương ứng). Do
đó, biến này được duy trì ở định dạng số trong suốt quá trình chạy mô hình.

Về biến \texttt{manufacturer}: Trong dữ liệu gốc, hãng ATI (đã bị AMD sáp nhập) xuất
hiện ở một số quan sát. Để đảm bảo tính nhất quán về thị phần hiện tại, nhóm đã tiến
hành gộp toàn bộ giá trị ``ATI'' vào nhóm ``AMD'' trong quá trình làm sạch dữ liệu.

% ------------------------------------------------------------
\subsection{Các bước xử lý dữ liệu}
% ------------------------------------------------------------

Quá trình xử lý dữ liệu được thực hiện theo trình tự sau, mỗi bước đều có lý do cụ thể
và được triển khai bằng ngôn ngữ R:

% ......................
\subsubsection*{Bước 1: Tải và khảo sát sơ bộ dữ liệu}
% ......................

\textbf{Kỹ thuật sử dụng:} Đọc file CSV bằng hàm \texttt{read.csv()} trong R, sau đó
dùng \texttt{str()}, \texttt{summary()} và \texttt{colSums(is.na(.))} để khảo sát cấu
trúc, kiểu dữ liệu và tỷ lệ giá trị khuyết của từng cột.

\textbf{Lý do:} Trước khi xử lý, cần nắm rõ ``hình dạng'' thực sự của dữ liệu -- số
dòng/cột, kiểu dữ liệu có đúng không (ví dụ: cột số bị đọc thành chuỗi do ký tự đặc
biệt), và cột nào bị thiếu nhiều đến mức không dùng được. Đây là bước bắt buộc trong
bất kỳ quy trình phân tích dữ liệu nào.

% ......................
\subsubsection*{Bước 2: Chọn lọc cột và ép kiểu dữ liệu}
% ......................

\textbf{Kỹ thuật sử dụng:} Subset dataframe để chỉ giữ lại 8 cột cần thiết. Dùng
\texttt{as.numeric()} để ép các cột số bị lưu sai kiểu, và parse ngày phát hành bằng
\texttt{as.integer(format(as.Date(release\_date), "\%Y"))} để trích xuất năm dưới dạng
số nguyên.

\textbf{Lý do:} Giảm chiều dữ liệu giúp tập trung vào những biến có ý nghĩa và giảm
nhiễu trong mô hình. Việc ép kiểu đúng là điều kiện tiên quyết để R thực hiện phép tính
số học và thống kê chính xác. Cụ thể, \texttt{release\_date} dạng chuỗi
``\texttt{05-Mar-2017}'' (ví dụ: ngày ra mắt của GeForce GTX 1080 Ti) cần được chuyển về
số nguyên \texttt{2017} lưu vào biến mới \texttt{release\_year} vì mô hình hồi quy cần
giá trị số liên tục, không phải chuỗi thời gian.

% ......................
\subsubsection*{Bước 3: Phát hiện và xử lý ngoại lệ (Outliers)}
% ......................

\textbf{Kỹ thuật sử dụng:} Sử dụng phương pháp IQR (Interquartile Range) -- quan sát
được coi là ngoại lệ nếu nằm ngoài khoảng $[Q_1 - 1{,}5 \times \mathrm{IQR},\ Q_3 +
1{,}5 \times \mathrm{IQR}]$. Áp dụng cho biến phụ thuộc \texttt{release\_price} và các
biến số liên tục có phân phối lệch mạnh (\texttt{tdp}, \texttt{core\_speed}). Bước này
bắt buộc phải thực hiện trước khi điền khuyết (imputation) nhằm tính toán các ranh giới
ngoại lệ (IQR bounds) dựa trên phân phối gốc nguyên thủy. Nếu điền khuyết hàng loạt
bằng trung vị trước, dải phân vị sẽ bị thu hẹp giả tạo, dẫn đến việc nhận diện sai
ngoại lệ.

\textbf{Kết quả phát hiện:}
(\textit{Lưu ý về trình tự thực thi: Để ranh giới phân vị IQR không bị thu hẹp giả tạo,
hàm \texttt{count\_outliers()} được chạy ngay tại Bước~3 trên toàn bộ tập dữ liệu thô
$N = 3.406$ trước khi tiến hành điền khuyết. Tuy nhiên, để phản ánh đúng đặc thù dữ
liệu đưa vào mô hình, các con số báo cáo dưới đây là số lượng ngoại lệ được thống kê lại
trên tập phân tích chính thức $N = 556$ quan sát, sau khi đã lọc bỏ các dòng khuyết giá
bán -- xem chi tiết quy trình tại Bước~4}):

\begin{itemize}[leftmargin=1.5em]
  \item \textbf{\texttt{release\_price}}: phát hiện \textbf{33 ngoại lệ} (chiếm 5,9\%
        dữ liệu), chủ yếu là các GPU chuyên dụng hoặc dòng flagship tản nhiệt nước (ví
        dụ: Quadro Plex 7000 giá 14.999 USD, Radeon R9 295X2 giá 1.499 USD), có giá vượt
        xa mặt bằng chung.

  \item \textbf{\texttt{tdp}}: phát hiện \textbf{27 ngoại lệ} (giảm từ con số \textbf{94
        ngoại lệ} phát hiện \textbf{trên tập thô}), tương ứng với các GPU công suất cao
        vượt ngưỡng giới hạn trên của công thức IQR ($> 310\,\mathrm{W}$).

  \item \textbf{\texttt{core\_speed}}: phát hiện \textbf{29 ngoại lệ} (giảm từ con số
        \textbf{83 ngoại lệ} phát hiện \textbf{trên tập thô}), chủ yếu là các card có
        xung nhịp bất thường so với phân phối chung.
\end{itemize}

\textbf{Quyết định xử lý:}

\begin{itemize}[leftmargin=1.5em]
  \item Giữ lại tất cả các ngoại lệ của cả 3 biến (giá, tdp, core\_speed). Toàn bộ
        ngoại lệ này đều \textbf{chỉ nằm ở phía trên giới hạn trên} (Upper Bound) của
        công thức IQR, không có giá trị âm vô lý.

  \item Sự xuất hiện của chúng tương ứng với các quan sát hợp lệ thuộc nhóm GPU flagship,
        các dòng card tản nhiệt nước (công suất/xung nhịp cực cao) hoặc card chuyên dụng
        dành cho máy trạm. Vì đây là đặc thù phân phối lệch phải tự nhiên của thị trường
        phần cứng (không phải lỗi nhập liệu), nhóm quyết định giữ nguyên tập dữ liệu
        $N = 556$ để phản ánh đúng thực tế.
\end{itemize}

\textbf{Lý do:} Ngoại lệ có thể làm hệ số hồi quy bị lệch đáng kể (đặc biệt với OLS --
Ordinary Least Squares vốn nhạy cảm với điểm dữ liệu xa). Tuy nhiên, không phải ngoại lệ
nào cũng là lỗi; cần phân biệt giữa lỗi nhập liệu và quan sát cực trị hợp lệ.

% ......................
\subsubsection*{Bước 4: Xử lý giá trị khuyết (Missing Values)}
% ......................

\textbf{Kỹ thuật sử dụng:} Chiến lược xử lý được phân tầng theo vai trò của biến:

\begin{itemize}[leftmargin=1.5em]
  \item \textbf{Biến phụ thuộc (\texttt{release\_price}):} Xóa toàn bộ dòng có giá trị
        NA vì không thể nội suy biến mục tiêu, điều này sẽ làm sai lệch mô hình.

  \item \textbf{Biến số liên tục (\texttt{tdp}, \texttt{memory\_size},
        \texttt{memory\_bus}):} Thay thế NA bằng trung vị (median) của từng cột. Trung vị
        được ưu tiên hơn trung bình vì dữ liệu GPU thường có phân phối lệch phải
        (skewed right) do sự xuất hiện của các dòng sản phẩm cao cấp. Đồng thời, vì nhóm
        đã quyết định giữ lại các ngoại lệ (flagship GPU) ở Bước~3, việc sử dụng trung vị
        càng thể hiện sự ưu việt vì nó là đại lượng thống kê bền vững (robust), không bị
        kéo lệch bởi các giá trị cực trị này khi dùng cho median imputation.

  \item \textbf{Biến phân loại (\texttt{manufacturer}, \texttt{memory\_type}):}
        (\textit{Lưu ý: Việc gộp nhãn ATI vào AMD được thực hiện trước bước này}).  Tiến
        hành thay thế NA bằng giá trị xuất hiện nhiều nhất (mode). Phương pháp này bảo
        toàn phân phối tần suất của biến phân loại tốt hơn so với xóa dòng.
\end{itemize}

\textbf{Thống kê giá trị khuyết thực tế từ bộ dữ liệu} (kết quả
\texttt{colSums(is.na(.))} sau bước chọn lọc 8 biến):

\begin{itemize}[leftmargin=1.5em]
  \item \textbf{\texttt{release\_price}} có tỷ lệ NA cao nhất (83,7\% -- 2.850/3.406
        quan sát), phản ánh thực tế nhiều GPU không công bố giá chính thức.

  \item \textbf{\texttt{Max\_Power} (\texttt{tdp})} thiếu 18,3\% (625 dòng);
        \textbf{\texttt{Memory} (\texttt{memory\_size})} thiếu 12,3\% (420 dòng);
        \textbf{\texttt{Memory\_Bus}} thiếu 1,8\% (62 dòng); \textbf{\texttt{Memory\_Type}}
        thiếu 1,6\% (56 dòng).

  \item \textbf{\texttt{Manufacturer}} không có giá trị khuyết (0\%). Riêng biến
        \textbf{\texttt{core\_speed}}, mặc dù dữ liệu thô ban đầu không có ô trống (0\%),
        nhưng chứa nhiều chuỗi ký tự không hợp lệ (ví dụ: ``\verb|\n- |''). Sau khi ép
        kiểu về dạng số (numeric), các giá trị không đo được này biến thành NA và được
        nhóm xử lý đồng nhất theo phương pháp điền trung vị.
\end{itemize}

Sau khi xóa các dòng có NA ở \texttt{release\_price} (biến phụ thuộc), tập dữ liệu phân
tích còn \textbf{556 quan sát hợp lệ}. Đối với các biến độc lập còn thiếu, chúng tôi
tiến hành imputation như trên.

\textbf{Nhận xét về tính đại diện của mẫu:} Việc loại bỏ hơn 83\% quan sát do thiếu giá
phát hành có thể dẫn đến nguy cơ sai lệch chọn mẫu (selection bias). Thực tế, các GPU
thiếu giá không phân bố ngẫu nhiên mà thường tập trung ở các dòng card cũ hoặc card đồ
họa tích hợp (điển hình là toàn bộ 254 quan sát của hãng Intel trong dữ liệu gốc đều rơi
vào trường hợp này và bị loại bỏ hoàn toàn). Đây là hạn chế cố hữu xuất phát từ đặc
điểm của bộ dữ liệu gốc, giá phát hành không được nhà sản xuất công bố chính thức cho
toàn bộ sản phẩm, dẫn đến tỷ lệ NA cao mang tính cấu trúc chứ không phải ngẫu nhiên
(Missing Not At Random --- MNAR). Vì vậy, nhóm ghi nhận đây là giới hạn của nghiên cứu;
các kết luận từ mô hình mang tính chất khám phá cấu trúc định giá của phân khúc GPU bán
lẻ và nên được diễn giải một cách thận trọng, không nên được ngoại suy một cách tuyệt đối
ra toàn bộ thị trường GPU.

% ......................
\subsubsection*{Bước 5: Mã hóa biến phân loại (Encoding)}
% ......................

\textbf{Kỹ thuật sử dụng:} Chuyển các biến phân loại \texttt{manufacturer} và
\texttt{memory\_type} thành \texttt{factor} trong R bằng \texttt{as.factor()}. Khi đưa
vào hàm \texttt{lm()} (hồi quy tuyến tính), R tự động tạo biến giả (dummy variables)
với một nhóm làm cơ sở (reference level).

\textbf{Lý do:} Các thuật toán hồi quy tuyến tính chỉ nhận đầu vào là số. Biến phân loại
như ``NVIDIA'', ``AMD'' không thể đưa trực tiếp vào phương trình, mà phải được mã hóa
thành các cột 0/1. R xử lý tự động việc này khi biến đã được khai báo là factor, giúp
đơn giản hóa code và tránh lỗi.

% ......................
\subsubsection*{Bước 6: Chuẩn hóa dữ liệu (Normalization / Standardization)}
% ......................

\textbf{Kỹ thuật sử dụng:} Chuẩn hóa z-score (standardization) cho tất cả các biến số
độc lập liên tục (\texttt{tdp}, \texttt{core\_speed}, \texttt{memory\_size},
\texttt{memory\_bus}):
\[
  z = \frac{x - \mu}{\sigma}
\]
được thực hiện bằng hàm \texttt{scale()} trong R.

\textbf{Lý do:} Các biến có đơn vị đo và biên độ khác nhau (TDP tính bằng Watt trong
khoảng $18 - 600$, \texttt{core\_speed} tính bằng MHz trong khoảng $550 - 1784\ldots$)
sẽ tạo ra hệ số hồi quy không so sánh được với nhau.

\textbf{Lưu ý:} Đối với biến \texttt{release\_year} (năm phát hành), sau khi parse, có
dạng số nguyên. Để tối ưu hóa khả năng diễn giải thực tế của mô hình hồi quy tuyến tính
(đo lường trực tiếp mức thay đổi giá trị GPU sau mỗi một năm trôi qua), nhóm quyết định
\textbf{giữ nguyên giá trị gốc} của biến này thay vì áp dụng chuẩn hóa Z-score như các
biến thông số kỹ thuật khác. Việc này giúp các hệ số mô hình mang ý nghĩa kinh tế học
trực quan hơn.

% ------------------------------------------------------------
\subsection{Xây dựng mô hình thống kê}
% ------------------------------------------------------------

\subsubsection{Mô hình hồi quy tuyến tính bội (Multiple Linear Regression)}

\textbf{Mục tiêu:} Xây dựng phương trình mô hình hóa giá phát hành GPU
(\texttt{release\_price}) từ sự tác động đồng thời của các thông số kỹ thuật cốt lõi.
Quá trình này bao gồm việc thử nghiệm và tinh chỉnh danh sách các biến độc lập ban đầu
(như TDP, xung nhịp, bộ nhớ, năm phát hành, hãng sản xuất) để loại bỏ các yếu tố gây
nhiễu hoặc đa cộng tuyến, từ đó xác định chính xác thông số nào có ảnh hưởng mạnh nhất
đến giá.

\textbf{Lý do chọn mô hình:} Trong bối cảnh nghiên cứu này, mục tiêu chính không phải là
dự đoán giá ngoại mẫu mà quan trọng hơn là hiểu cấu trúc định giá và lượng hóa tác động
riêng phần của từng thông số. Hồi quy tuyến tính bội được ưu tiên vì tính diễn giải cao
(interpretability), mỗi hệ số hồi quy trực tiếp thể hiện mức thay đổi về giá khi một
thông số thay đổi \textbf{một độ lệch chuẩn (1~SD)} (sau khi đã kiểm soát các yếu tố
khác). Điều này phù hợp với nhu cầu của nhà sản xuất và nhà phân tích muốn biết chính xác
yếu tố kỹ thuật nào đóng góp nhiều nhất vào giá thành. Các mô hình machine learning phức
tạp (như random forest, neural network) tuy có thể đạt độ phù hợp dữ liệu (fit) cao hơn
nhưng lại khó giải thích được cơ chế tác động, do đó không được lựa chọn trong phạm vi
bài báo cáo này.

\textbf{Giới hạn của mô hình:} Mặc dù hồi quy tuyến tính cung cấp cách diễn giải trực
quan, nhưng giả định tuyến tính có thể không nắm bắt được các quan hệ phi tuyến (ví dụ:
hiệu suất tăng phi tuyến theo băng thông bộ nhớ) hoặc tương tác giữa các biến (ví dụ:
tác động của TDP lên giá có thể phụ thuộc vào hãng sản xuất). Việc kiểm tra hiệu ứng phi
tuyến sâu hơn (như kiểm định Ramsey RESET) được đề xuất như một hướng mở rộng cho nghiên
cứu tiếp theo.

\textbf{Kỹ thuật triển khai:} Sử dụng hàm \texttt{lm()} trong R để ước lượng các hệ số
tác động.

\textbf{Đánh giá và tinh chỉnh mô hình:} Sau khi chạy mô hình sơ bộ, nhóm tiến hành các
bước kiểm định bắt buộc:

\begin{itemize}[leftmargin=1.5em]
  \item \textbf{Mức độ giải thích của mô hình} được đo bằng $R^2$ hiệu chỉnh (Adjusted
        $R^2$) để biết tỷ lệ phần trăm của phương sai của giá GPU được giải thích bởi các
        thông số kỹ thuật.

  \item \textbf{Đa cộng tuyến} được đánh giá bằng chỉ số VIF. Trong thực hành thống kê
        khắt khe, VIF $> 5$ đã bắt đầu cho thấy các biến độc lập có tương quan quá chặt
        chẽ, gây nhiễu và làm giảm độ chính xác của ước lượng hệ số. Dựa trên tiêu chuẩn
        này kết hợp với nguyên tắc tiết kiệm mô hình (Occam's Razor), nhóm chủ động loại
        bỏ biến không đạt yêu cầu (cụ thể là biến \texttt{memory\_type} với VIF = 8.12)
        để đảm bảo tính độc lập, tính vững (robustness) và độ tin cậy tối đa cho các hệ
        số hồi quy còn lại.

  \item \textbf{Phương sai sai số không đổi} được kiểm định bằng Breusch-Pagan. Nếu vi
        phạm, các ước lượng hệ số vẫn không chệch nhưng sai số chuẩn sẽ bị ảnh hưởng,
        làm sai lệch kết luận về ý nghĩa thống kê. Đây là hạn chế của OLS truyền thống;
        nhóm xử lý bằng phép biến đổi log cho biến $Y$ thay vì dùng phương sai sai số
        chuẩn vững (Robust SE).

  \item \textbf{Phân phối chuẩn của phần dư} được kiểm định bằng Shapiro-Wilk. Do đặc
        thù dữ liệu giá cả thường gây ra phần dư lệch biên, nhóm chủ động áp dụng phép
        biến đổi logarit cho biến phụ thuộc thành $\log(\texttt{release\_price})$. Kỹ
        thuật này giúp vừa cải thiện tiệm cận chuẩn cho phân phối phần dư, vừa ổn định
        phương sai sai số.
\end{itemize}

\subsubsection{Phân tích phương sai (Welch's ANOVA)}

\textbf{Mục tiêu:} Kiểm định xem có sự chênh lệch có ý nghĩa thống kê về giá phát hành
trung bình giữa hai nhóm nhà sản xuất (NVIDIA và AMD) hay không. Phân tích này giúp trả
lời độc lập câu hỏi liệu thương hiệu có phải là yếu tố tạo ra khoảng cách giá trước khi
đưa vào mô hình đa biến.

\textbf{Từ góc độ kinh tế học:} Nếu kết quả cho thấy sự khác biệt có ý nghĩa giữa hai
hãng, đó là minh chứng thống kê cho chiến lược định giá thương hiệu (ví dụ: xác nhận giả
thuyết NVIDIA áp dụng mức giá trung bình cao hơn AMD). Điều này cung cấp insight quan
trọng cho cả người tiêu dùng và nhà định giá.

\textbf{Kiểm định tiền điều kiện:} Do cỡ mẫu lớn (NVIDIA $n=284$, AMD $n=272$), điều
kiện phân phối chuẩn nội nhóm được đảm bảo một cách mạnh mẽ theo Định lý Giới hạn Trung
tâm (CLT). Nhóm tiến hành kiểm tra tính đồng nhất của phương sai bằng Levene's test. Do
đặc thù phân khúc sản phẩm của hai hãng khác nhau, phương sai giá dự kiến sẽ không đồng
nhất.

\textbf{Kỹ thuật triển khai:} Do chỉ có 2 nhóm quan sát và kết quả kiểm định Levene's
Test trước đó đã xác nhận sự vi phạm giả định đồng nhất phương sai, nhóm quyết định sử
dụng kiểm định mạnh (robust) là \textbf{Welch's ANOVA} (hàm \texttt{oneway.test()} trong
R). Mức ý nghĩa được chọn là $\alpha = 0.05$. (\textit{Lưu ý: Vì tập phân tích chỉ còn
2 nhóm, Welch's ANOVA về mặt toán học cho kết quả hoàn toàn tương đương với phép kiểm
định Welch's T-test, do đó không yêu cầu thực hiện thêm kiểm định hậu kiểm Post-hoc.})

% ------------------------------------------------------------
\subsection{Sơ đồ tổng quan quy trình xử lý}
% ------------------------------------------------------------

Quy trình phân tích dữ liệu của nhóm được tóm tắt trực quan qua sơ đồ khối dưới đây:

\begin{center}
  \textit{Hình 1. Sơ đồ quy trình xử lý và phân tích dữ liệu}
\end{center}

% ============================================================
\section*{Tài liệu tham khảo}
\addcontentsline{toc}{section}{Tài liệu tham khảo}
% ============================================================

\begin{enumerate}[leftmargin=1.5em, label={[\arabic*]}]

  \item Iliasse Kkaf. (2017). \textit{Computer Parts (CPUs and GPUs)}. Kaggle. Truy cập
        tại: \url{https://www.kaggle.com/datasets/iliassekkaf/computerparts}. Truy cập lần
        cuối: tháng 4/2026.

  \item Montgomery, D.\,C., Peck, E.\,A., \& Vining, G.\,G. (2012). \textit{Introduction
        to Linear Regression Analysis} (5th ed.). John Wiley \& Sons.

  \item R Core Team. (2024). \textit{R: A Language and Environment for Statistical
        Computing}. R Foundation for Statistical Computing, Vienna, Austria.
        \url{https://www.R-project.org/}.

  \item Tukey, J.\,W. (1977). \textit{Exploratory Data Analysis}. Addison-Wesley.

  \item Breusch, T.\,S., \& Pagan, A.\,R. (1979). \textit{Econometrica, 47}(5),
        1287--1294.

  \item Shapiro, S.\,S., \& Wilk, M.\,B. (1965). \textit{Biometrika, 52}(3--4),
        591--611.

  \item Levene, H. (1960). In \textit{Contributions to Probability and Statistics}.
        Stanford University Press.

  \item Welch, B.\,L. (1951). \textit{Biometrika, 38}(3--4), 330--336.

  \item Ramsey, J.\,B. (1969). \textit{Journal of the Royal Statistical Society,
        31}(2), 350--371.

  \item O'Brien, R.\,M. (2007). \textit{Quality \& Quantity, 41}(5), 673--690.

\end{enumerate}

\end{document}